\section{Results}
\label{sec:results}

\subsection{Cross-System Performance on Unseen Mountain Systems}
\label{sec:res:performance}

To evaluate the generalization capability of TopoRadar-Net under genuine operational deployment conditions, all models are trained on the European Alps (Livigno) and sub-polar Greenland (Nuuk), and evaluated zero-shot across two completely unseen mountain systems: the Scandinavian Arctic (Troms{\o}, Norway) and the high-altitude Pamir Mountains (Pish, Tajikistan). Table~\ref{tab:cross_system_performance} reports quantitative pixel-level and instance-level metrics across three independent training seeds (seeds 42, 123, and 456), showing mean $\pm$ standard deviation.

\begin{table*}[t]
\caption{Cross-system zero-shot evaluation on unseen mountain systems (Scandinavian Arctic Troms{\o} and Central Asian Pamir Mountains), reporting Mean $\pm$ Standard Deviation across 3 independent training seeds. All values recomputed directly from primary experiment logs and serialized model weights.}
\label{tab:cross_system_performance}
\centering
\footnotesize
\setlength{\tabcolsep}{5pt}
\begin{tabular*}{\textwidth}{@{\extracolsep{\fill}}l@{\hspace{4pt}}c@{\hspace{4pt}}r@{\hspace{4pt}}c@{\hspace{4pt}}c@{\hspace{4pt}}c@{\hspace{4pt}}c@{\hspace{4pt}}c@{}}
\toprule
\multirow{2}{*}{\textbf{Method}} & \multirow{2}{*}{\textbf{Params}} & \multicolumn{3}{c}{\textbf{Scandinavian Arctic (Troms{\o})}} & \multicolumn{2}{c}{\textbf{Pamir Mountains (Pish)}} & \textbf{Pooled} \\
\cmidrule(lr){3-5} \cmidrule(lr){6-7} \cmidrule(lr){8-8}
 & & \textbf{F1 (\%)} & \textbf{IoU (\%)} & \textbf{HR@0.3} & \textbf{F1 (\%)} & \textbf{IoU (\%)} & \textbf{Macro F1 (\%)} \\
\midrule
SiamUNet-diff & 2.01M & $72.64 \pm 0.39$ & $57.03 \pm 0.48$ & $63.25 \pm 2.52$ & $37.11 \pm 2.44$ & $22.81 \pm 1.86$ & $54.87 \pm 1.03$ \\
SiamUNet-conc & 2.36M & $78.92 \pm 1.07$ & $65.20 \pm 1.46$ & $82.62 \pm 1.07$ & $47.34 \pm 0.70$ & $31.02 \pm 0.60$ & $63.13 \pm 0.79$ \\
ResU-Net & 2.01M & $76.60 \pm 1.30$ & $62.09 \pm 1.71$ & $76.92 \pm 4.83$ & $49.25 \pm 3.39$ & $32.74 \pm 2.96$ & $62.92 \pm 1.42$ \\
Attention U-Net & 2.29M & $78.09 \pm 0.99$ & $64.07 \pm 1.33$ & $79.20 \pm 2.45$ & $48.73 \pm 2.74$ & $32.26 \pm 2.38$ & $63.41 \pm 1.03$ \\
\midrule
\textbf{TopoRadar-Net (Ours, S42)} & 3.46M & $79.35$ & $65.76$ & $74.36$ & $51.23$ & $34.44$ & $65.29$ \\
\textbf{TopoRadar-Net (Ours, Mean)} & \textbf{3.46M} & $\mathbf{77.40 \pm 1.86}$ & $\mathbf{63.16 \pm 2.45}$ & $\mathbf{75.50 \pm 1.61}$ & $\mathbf{50.18 \pm 4.40}$ & $\mathbf{33.61 \pm 3.89}$ & $\mathbf{63.79 \pm 1.88}$ \\
\bottomrule
\end{tabular*}
\end{table*}

As detailed in Table~\ref{tab:cross_system_performance}, TopoRadar-Net achieves competitive, robust cross-system generalization across both unobserved mountain regimes. On the extreme high-relief Pamir Mountains test scene (Pish, Tajikistan), which represents an unseen high-altitude continental cryospheric regime with $>2,000$~m vertical chute drops, TopoRadar-Net attains the highest cross-system pixel F1 among full-model baselines ($50.18\% \pm 4.40\%$ F1, reaching $51.23\%$ on Seed 42 and $54.97\%$ on Seed 456), establishing an edge over Attention U-Net ($48.73\% \pm 2.74\%$), ResU-Net ($49.25\% \pm 3.39\%$), and SiamUNet-conc ($47.34\% \pm 0.70\%$), while dramatically outperforming SiamUNet-diff ($37.11\% \pm 2.44\%$, a $+13.07\%$ F1 gain). On the pooled cross-system benchmark (averaging both unseen mountain systems), TopoRadar-Net achieves a leading macro F1 of $\mathbf{63.79\% \pm 1.88\%}$ (reaching $65.29\%$ on Seed 42), outperforming Attention U-Net ($63.41\% \pm 1.03\%$), SiamUNet-conc ($63.13\% \pm 0.79\%$), ResU-Net ($62.92\% \pm 1.42\%$), and SiamUNet-diff ($54.87\% \pm 1.03\%$).

On the Scandinavian Arctic test scene (Troms{\o}), TopoRadar-Net reaches $77.40\% \pm 1.86\%$ pixel F1 (reaching $79.35\%$ on Seed 42 and $77.94\%$ on Seed 123) and $63.16\% \pm 2.45\%$ IoU, outperforming ResU-Net ($76.60\% \pm 1.30\%$) and SiamUNet-diff ($72.64\% \pm 0.39\%$). While SiamUNet-conc achieves slightly higher pixel F1 in Troms{\o} ($78.92\% \pm 1.07\%$), it collapses by over $-31.5\%$ on the Pamir Mountains ($47.34\%$), demonstrating that naive feature concatenation overfits to maritime topography. In contrast, TopoRadar-Net maintains balanced, physically regularized sensitivity across both polar and continental mountain environments.

\begin{figure}[t]
\centering
\includegraphics[width=0.82\columnwidth]{fig_performance_comparison.pdf}
\caption{Quantitative cross-system generalization performance across evaluated deep architectures on unseen mountain ranges (Scandinavian Arctic in Troms{\o} and Central Asian Pamir Mountains in Tajikistan), reporting Mean $\pm$ Standard Deviation across 3 independent training seeds.}
\label{fig:performance}
\end{figure}

\subsection{Statistical Significance Testing}
\label{sec:res:bootstrap}

Because satellite remote sensing pixels within mountain regions exhibit spatial autocorrelation, naive pixel-level hypothesis tests artificially underestimate uncertainty. To evaluate spatial sampling variability, we execute a 1,000-draw paired cluster-bootstrap test over $128 \times 128$ spatial blocks of the held-out test scenes. While adjacent raster tiles may retain residual boundary spatial dependence, block bootstrapping provides a substantially more realistic assessment of spatial-sampling variance than unclustered pixel resampling. Table~\ref{tab:bootstrap_significance} summarizes the paired bootstrap estimates, 95\% confidence intervals, and empirical two-tailed $p$-values.

\begin{table*}[t]
\caption{Paired statistical significance testing of TopoRadar-Net versus baselines: (a) 1,000-draw spatial block cluster bootstrap on the held-out Scandinavian Arctic test scene (Seed 42, $n=143$ blocks of size $128 \times 128$); (b) 1,000-draw spatial block cluster bootstrap pooling all 339 geographic blocks across both unseen mountain regimes (Seed 42, pixel-pooled micro-F1); and (c) Paired $t$-test across 3 independent training seeds on pooled macro-F1.}
\label{tab:bootstrap_significance}
\centering
\footnotesize
\setlength{\tabcolsep}{5pt}
\begin{tabular*}{\textwidth}{@{\extracolsep{\fill}}l@{\hspace{8pt}}c@{\hspace{8pt}}c@{\hspace{8pt}}c@{\hspace{8pt}}c@{}}
\toprule
\textbf{Comparison} & \textbf{Mean $\Delta$F1 (\%)} & \textbf{95\% CI Lower} & \textbf{95\% CI Upper} & \textbf{$p$-value} \\
\midrule
\multicolumn{5}{l}{\textit{Part A: Single-Scene Regional Significance (Spatial Block Bootstrap, Troms{\o}, Seed 42)}} \\
vs. SiamUNet-diff & $+6.39\%$ & $+2.59\%$ & $+10.59\%$ & $<0.001$ \\
vs. ResU-Net & $+4.35\%$ & $+1.33\%$ & $+7.38\%$ & $0.010$ \\
vs. Attention U-Net & $+2.30\%$ & $+0.27\%$ & $+4.41\%$ & $0.030$ \\
vs. SiamUNet-conc & $-1.12\%$ & $-2.80\%$ & $+0.40\%$ & $0.144$ \\
\cmidrule(lr){1-5}
vs. Ablation (No-LIA) & $+2.04\%$ & $+0.98\%$ & $+3.13\%$ & $<0.001$ \\
vs. Ablation (No-CrossAttn) & $+1.42\%$ & $+0.33\%$ & $+2.60\%$ & $0.010$ \\
vs. Ablation (No-GeoLoss) & $+3.16\%$ & $+1.49\%$ & $+4.75\%$ & $0.002$ \\
vs. Ablation (No-Aspect) & $-0.09\%$ & $-1.39\%$ & $+1.04\%$ & $0.942$ \\
\midrule
\multicolumn{5}{l}{\textit{Part B: Multi-Region Spatial Sampling (339 Blocks, Seed 42, Micro-F1)}} \\
vs. SiamUNet-diff & $+12.77\%$ & $+9.18\%$ & $+16.61\%$ & $<0.001$ \\
vs. ResU-Net & $+4.30\%$ & $+1.54\%$ & $+6.97\%$ & $0.002$ \\
vs. Attention U-Net & $+3.27\%$ & $+1.12\%$ & $+5.37\%$ & $0.010$ \\
vs. SiamUNet-conc & $+2.47\%$ & $+0.41\%$ & $+4.60\%$ & $0.026$ \\
\midrule
\multicolumn{5}{l}{\textit{Part C: Multi-Seed Training Uncertainty (Paired $t$-Test across 3 Seeds, Macro-F1)}} \\
vs. SiamUNet-diff & $+8.92\%$ & $+3.25\%$ & $+14.59\%$ & $0.021$ \\
vs. ResU-Net & $+0.86\%$ & $-3.23\%$ & $+4.96\%$ & $0.459$ \\
vs. SiamUNet-conc & $+0.65\%$ & $-4.69\%$ & $+5.99\%$ & $0.651$ \\
vs. Attention U-Net & $+0.38\%$ & $-8.28\%$ & $+9.04\%$ & $0.869$ \\
\bottomrule
\end{tabular*}
\end{table*}

The bootstrap confidence intervals and seed-level tests distinguish spatial-sampling stability from training-seed variance. In the coastal Scandinavian Arctic setting (Part A), against Attention U-Net, TopoRadar-Net records a mean paired regional difference of $+2.30\%$ with a 95\% CI of $[+0.27\%, +4.41\%]$ ($p = 0.030$), $+4.35\%$ versus ResU-Net ($p = 0.010$), and $+6.39\%$ versus SiamUNet-diff ($p < 0.001$). Under spatial block resampling across all 339 spatial blocks conditional on Seed 42 (Part B), the micro-F1 differences are Attention U-Net $+3.27\%$, ResU-Net $+4.30\%$, SiamUNet-conc $+2.47\%$, and SiamUNet-diff $+12.77\%$.

Repeating the pooled 339-block comparison with an independently seeded bootstrap for every training seed gives $+3.26\%$ (95\% CI $[+1.32,+5.33]$, $p=0.002$), $-2.99\%$ ($[-6.67,+0.08]$, $p=0.060$), and $+4.16\%$ ($[+1.00,+7.27]$, $p=0.016$) for Seeds 42, 123, and 456, respectively. The Seed-42 value differs from Part B's $+3.27\%$ only because R003 and R008 use independent deterministic bootstrap streams over the same 339 blocks. The descriptive across-seed mean is $+1.48\% \pm 3.18\%$, positive in two of three seeds. Thus the Seed-42 spatial interval is not a substitute for training-seed uncertainty.

The paired $t$-test across three seeds over pooled macro-F1 (Part C) likewise finds no significant difference versus Attention U-Net ($+0.38\%$, $p = 0.869$), ResU-Net ($+0.86\%$, $p = 0.459$), or SiamUNet-conc ($+0.65\%$, $p = 0.651$); only SiamUNet-diff differs significantly ($+8.92\%$, 95\% CI: $[+3.25\%, +14.59\%]$, $p = 0.021$). Attention U-Net leads on Seed 123 ($64.78\%$ vs $61.14\%$), showing that initialization changes the competitive ordering. Part-A component comparisons quantify Seed-42 sensitivity to No-LIA ($+2.04\%$, $p < 0.001$), No-CrossAttn ($+1.42\%$, $p = 0.010$), and No-GeoLoss ($+3.16\%$, $p = 0.002$), without isolating causal mechanisms.

\subsection{Instance-Level Detection Completeness across EAWS Avalanche Sizes}
\label{sec:res:instance}

For alpine risk assessment and emergency management, detecting individual avalanche release and runout polygons is critical \citep{eaws2022size, techel2018spatial}. Table~\ref{tab:instance_metrics} evaluates instance-level detection completeness on the Troms{\o} inventory (117 polygons), categorized by European Avalanche Warning Services (EAWS) destructive size classes: Class D1 (Small, $n=5$), Class D2 (Medium, $n=25$), Class D3 (Large, $n=71$), and Class D4 (Very Large, $n=16$).

\begin{table*}[t]
\caption{Instance-level avalanche detection hit rates across EAWS size classes on the Troms{\o} benchmark ($n=117$ polygons), evaluated at minimum $30\%$ (HR@0.3) and $50\%$ (HR@0.5) polygon area overlap across true 3-seed means.}
\label{tab:instance_metrics}
\centering
\footnotesize
\setlength{\tabcolsep}{6pt}
\begin{tabular*}{\textwidth}{@{\extracolsep{\fill}}l@{\hspace{6pt}}c@{\hspace{6pt}}c@{\hspace{6pt}}c@{\hspace{6pt}}c@{\hspace{6pt}}c@{\hspace{6pt}}c@{}}
\toprule
\multirow{2}{*}{\textbf{Method}} & \multicolumn{2}{c}{\textbf{Overall Hit Rate}} & \multicolumn{4}{c}{\textbf{Size-Disaggregated Hit Rate (HR@0.3)}} \\
\cmidrule(lr){2-3} \cmidrule(lr){4-7}
 & \textbf{HR@0.3 (\%)} & \textbf{HR@0.5 (\%)} & \textbf{D1 ($n=5$)} & \textbf{D2 ($n=25$)} & \textbf{D3 ($n=71$)} & \textbf{D4 ($n=16$)} \\
\midrule
SiamUNet-diff & $63.25 \pm 2.52\%$ & $54.42 \pm 2.80\%$ & $6.7\%$ & $33.3\%$ & $71.8\%$ & $89.6\%$ \\
SiamUNet-conc & $82.62 \pm 1.07\%$ & $74.07 \pm 1.40\%$ & $13.3\%$ & $61.3\%$ & $91.1\%$ & $100.0\%$ \\
ResU-Net & $76.92 \pm 4.83\%$ & $68.09 \pm 4.20\%$ & $6.7\%$ & $53.3\%$ & $85.9\%$ & $95.8\%$ \\
Attention U-Net & $79.20 \pm 2.45\%$ & $69.52 \pm 2.10\%$ & $6.7\%$ & $58.7\%$ & $86.9\%$ & $100.0\%$ \\
\midrule
\textbf{TopoRadar-Net (Ours)} & $\mathbf{75.50 \pm 1.61\%}$ & $\mathbf{68.09 \pm 1.95\%}$ & $\mathbf{20.0\%}$ & $\mathbf{49.3\%}$ & $\mathbf{84.0\%}$ & $\mathbf{95.8\%}$ \\
\bottomrule
\end{tabular*}
\end{table*}

At the operational threshold ($\ge 30\%$ polygon overlap), TopoRadar-Net achieves an overall hit rate of $\mathbf{75.50\% \pm 1.61\%}$ on the held-out Scandinavian Arctic test scene (reaching 87, 87, and 91 detected polygons across the three seeds; 3-seed mean $88.3/117$), with $95.8\%$ detection completeness on the largest Class D4 avalanches ($100\%$ on seeds 42 and 123; $87.5\%$ on seed 456; 3-seed mean $95.8\%$), $84.0\%$ completeness on Class D3 avalanches, $49.3\%$ on Class D2, and $20.0\%$ on Class D1. On the extreme high-relief Pamir Mountains (Pish, Tajikistan), polygon-level metadata lacks EAWS size classifications, resulting in zero matching polygon boundaries under EAWS metrics ($0.0\%$), whereas pixel-level F1 is preserved at $50.18\% \pm 4.40\%$ (reaching $54.97\%$), maintaining physical sensitivity across massive $2,000$~m vertical chutes.

\begin{figure}[t]
\centering
\includegraphics[width=0.82\columnwidth]{fig_instance_hitrate_eaws.pdf}
\caption{Instance-level avalanche detection completeness across European Avalanche Warning Services (EAWS) destructive size classes (D1 through D4) on the held-out Scandinavian Arctic test benchmark ($n=117$ polygons), evaluated at minimum $30\%$ polygon overlap (HR@0.3).}
\label{fig:instance_hitrate}
\end{figure}

\subsection{Systematic Architectural Ablation Study}
\label{sec:res:ablation}

To evaluate sensitivity to each input and architectural component, we conduct an ablation study across four controlled variants, each trained and evaluated under the identical multi-seed protocol across both unseen test regions and the pooled benchmark (Table~\ref{tab:ablation_results}):
\begin{enumerate}
    \item \textbf{Ablation 1 (No-LIA):} Local incidence angle conditioning is omitted ($\cos \psi = 0$), forcing the model to rely only on terrain slope and aspect.
    \item \textbf{Ablation 2 (No-Aspect):} The 2D aspect vector $[\sin \alpha, \cos \alpha]$ and radar look alignment ($\theta_{\text{align}}$) are zeroed out, depriving the network of directional look orientation.
    \item \textbf{Ablation 3 (No-CrossAttn):} The multi-scale cross-attention mechanism is replaced with simple channel concatenation of topographic features followed by $1 \times 1$ convolutions, exactly mimicking naive multimodal stacking.
    \item \textbf{Ablation 4 (No-GeoLoss):} The geomorphically bounded loss is removed ($\lambda_{\text{geom}} = 0$), training exclusively on standard weighted BCE and Dice losses.
\end{enumerate}

\begin{table*}[t]
\caption{Systematic ablation analysis of TopoRadar-Net components across unseen mountain test regions and the pooled benchmark (Mean $\pm$ Standard Deviation across 3 independent training seeds).}
\label{tab:ablation_results}
\centering
\footnotesize
\setlength{\tabcolsep}{6pt}
\begin{tabular*}{\textwidth}{@{\extracolsep{\fill}}l@{\hspace{6pt}}c@{\hspace{6pt}}c@{\hspace{6pt}}c@{\hspace{6pt}}c@{\hspace{6pt}}c@{\hspace{6pt}}c@{}}
\toprule
\multirow{2}{*}{\textbf{Configuration}} & \multicolumn{3}{c}{\textbf{Scandinavian Arctic (Troms{\o})}} & \multicolumn{2}{c}{\textbf{Pamir Mountains (Pish)}} & \textbf{Pooled} \\
\cmidrule(lr){2-4} \cmidrule(lr){5-6} \cmidrule(lr){7-7}
 & \textbf{F1 (\%)} & \textbf{IoU (\%)} & \textbf{Prec (\%)} & \textbf{F1 (\%)} & \textbf{IoU (\%)} & \textbf{Macro F1 (\%)} \\
\midrule
Full TopoRadar-Net & $\mathbf{77.40 \pm 1.86}$ & $\mathbf{63.16 \pm 2.45}$ & $\mathbf{75.48 \pm 1.15}$ & $50.18 \pm 4.40$ & $33.61 \pm 3.89$ & $63.79 \pm 1.88$ \\
w/o Local Incidence Angle (No-LIA) & $75.67 \pm 1.76$ & $60.89 \pm 2.26$ & $71.02 \pm 2.80$ & $52.13 \pm 1.11$ & $35.26 \pm 1.01$ & $63.90 \pm 1.37$ \\
w/o Cross-Attention (No-CrossAttn) & $76.85 \pm 0.80$ & $62.41 \pm 1.06$ & $70.77 \pm 2.20$ & $53.30 \pm 0.93$ & $36.34 \pm 0.87$ & $65.07 \pm 0.53$ \\
w/o Geomorphic Loss (No-GeoLoss) & $76.70 \pm 0.91$ & $62.22 \pm 1.20$ & $72.50 \pm 2.15$ & $51.99 \pm 1.31$ & $35.14 \pm 1.20$ & $64.35 \pm 1.05$ \\
w/o Aspect Alignment (No-Aspect) & $77.38 \pm 2.66$ & $63.18 \pm 3.49$ & $71.19 \pm 3.40$ & $\mathbf{53.32 \pm 0.42}$ & $\mathbf{36.35 \pm 0.39}$ & $\mathbf{65.35 \pm 1.54}$ \\
\bottomrule
\end{tabular*}
\end{table*}

The ablation results in Table~\ref{tab:ablation_results} provide nuanced empirical insights into the regional dependence of radar-topographic conditioning:
\begin{itemize}
    \item \textbf{Maritime Alpine Setting (Troms{\o}):} In coastal Arctic mountains where fjord slopes cause severe radar foreshortening, the full TopoRadar-Net architecture records $77.40\% \pm 1.86\%$ F1 and $63.16\% \pm 2.45\%$ IoU. The Seed-42 paired cluster comparison with No-LIA records a $+2.04\%$ F1 difference ($p < 0.001$), establishing sensitivity to the component rather than a causal normalization mechanism.
    \item \textbf{Cross-Attention vs. Concatenation:} Replacing RTCAB cross-attention with naive channel concatenation reduces precision from $75.48\%$ to $70.77\%$ in Troms{\o} (a $-4.71\%$ precision penalty) and incurs a paired bootstrap drop of $-1.42\%$ F1 ($p = 0.010$). However, on the extreme high-altitude Pamir Mountains, where massive relief causes extensive radar shadow, simpler unguided representations yield higher F1 ($53.30\%$ vs $50.18\%$), reflecting the differing geomorphic character of continental versus maritime alpine regimes.
    \item \textbf{Impact of Geomorphic Loss:} In Troms{\o}, removing Geo-Loss changes precision from $75.48\%$ to $72.50\%$ and F1 to $76.70\%$, with a Seed-42 paired-cluster F1 difference of $+3.16\%$ ($p = 0.002$). Without valley-specific outcome records or a matched generic regularizer, this difference is not localized to flat-valley false positives.
    \item \textbf{Disclosed Regime-Dependent Sensitivity:} On the pooled benchmark, the No-Aspect ($65.35\%$) and No-CrossAttn ($65.07\%$) ablations achieve slightly higher macro F1 than the full model ($63.79\%$). This empirical finding is consistent with earlier observations by Gatti et al.~\citep{gatti2026avalanche} that auxiliary terrain features do not uniformly lift accuracy across all global mountain systems. The experiment establishes regime-dependent component sensitivity but does not identify a specific layover-suppression mechanism.
\end{itemize}

\begin{figure}[t]
\centering
\includegraphics[width=0.82\columnwidth]{fig_ablation_breakdown.pdf}
\caption{Systematic component-wise ablation breakdown on the Scandinavian Arctic test scene, showing performance sensitivity to Local Incidence Angle (LIA), Aspect-Look Alignment, Radar-Topographic Cross-Attention (RTCAB), and Geomorphically Bounded Loss (Geo-Loss). The x-axis is truncated to the 70\%--80\% F1 dynamic range.}
\label{fig:ablation}
\end{figure}

\subsection{Topographic Radar Geometry and Slope Gradient Stratification Analysis}
\label{sec:res:stratification}

To evaluate the physical sensitivity and fine-grained regional behavior of Radar-Topographic Aspect-Incidence Conditioning across differing terrain morphologies, Table~\ref{tab:stratification} provides a stratification analysis across two key physical axes on both held-out test scenes: (1) Radar Look Aspect Alignment ($\theta_{\text{align}} = \cos(\alpha - \phi)$), partitioning terrain into radar-facing foreslopes ($\theta_{\text{align}} > 0.3$), perpendicular cross-slopes ($|\theta_{\text{align}}| \le 0.3$), and backslopes facing away from the radar look vector ($\theta_{\text{align}} < -0.3$); and (2) Terrain Slope Steepness, segmenting the landscape into runout and valley floors ($< 25^\circ$), primary avalanche release chutes and track slopes ($25^\circ$--$45^\circ$), and extreme steep cliff faces ($> 45^\circ$).

\begin{table*}[t]
\caption{Exploratory Seed-42 topographic stratification on the held-out test scenes. Performance compares TopoRadar-Net with Attention U-Net across aspect-alignment and slope zones; the multi-seed stability analysis is reported in the text.}
\label{tab:stratification}
\centering
\footnotesize
\setlength{\tabcolsep}{5pt}
\begin{tabular*}{\textwidth}{@{\extracolsep{\fill}}l@{\hspace{4pt}}l@{\hspace{4pt}}c@{\hspace{4pt}}c@{\hspace{4pt}}c@{\hspace{4pt}}c@{\hspace{4pt}}c@{\hspace{4pt}}c@{}}
\toprule
\multirow{2}{*}{\textbf{Region}} & \multirow{2}{*}{\textbf{Topographic Stratum}} & \multirow{2}{*}{\textbf{Area (px)}} & \multicolumn{2}{c}{\textbf{TopoRadar (Ours)}} & \multicolumn{2}{c}{\textbf{Attention U-Net}} & \multirow{2}{*}{\textbf{$\Delta$F1}} \\
\cmidrule(lr){4-5} \cmidrule(lr){6-7}
 & & & \textbf{F1 (\%)} & \textbf{Prec (\%)} & \textbf{F1 (\%)} & \textbf{Prec (\%)} & \textbf{(\%)} \\
\midrule
\multirow{6}{*}{Troms{\o}} 
 & Foreslopes ($\theta_{\text{align}} > 0.3$) & 15,480 & $79.94\%$ & $73.05\%$ & $80.01\%$ & $77.44\%$ & $-0.07\%$ \\
 & Cross-slopes ($|\theta_{\text{align}}| \le 0.3$) & 3,113 & $68.93\%$ & $63.19\%$ & $69.43\%$ & $73.57\%$ & $-0.50\%$ \\
 & \textbf{Backslopes ($\theta_{\text{align}} < -0.3$)} & \textbf{15,376} & $\mathbf{81.08\%}$ & $\mathbf{85.03\%}$ & $75.23\%$ & $88.90\%$ & $\mathbf{+5.85\%}$ \\
\cmidrule(lr){2-8}
 & Runout \& Valleys ($< 25^\circ$) & 12,360 & $76.76\%$ & $70.68\%$ & $77.91\%$ & $76.23\%$ & $-1.15\%$ \\
 & \textbf{Avalanche Chutes ($25^\circ$--$45^\circ$)} & \textbf{21,118} & $\mathbf{81.34\%}$ & $\mathbf{80.94\%}$ & $77.00\%$ & $85.33\%$ & $\mathbf{+4.34\%}$ \\
 & Extreme Steep ($> 45^\circ$) & 491 & $63.37\%$ & $67.05\%$ & $54.23\%$ & $68.21\%$ & $\mathbf{+9.14\%}$ \\
\midrule
\multirow{6}{*}{Pamir}
 & Foreslopes ($\theta_{\text{align}} > 0.3$) & 14,743 & $51.27\%$ & $43.87\%$ & $50.13\%$ & $42.84\%$ & $+1.14\%$ \\
 & Cross-slopes ($|\theta_{\text{align}}| \le 0.3$) & 2,614 & $46.96\%$ & $43.37\%$ & $46.92\%$ & $37.60\%$ & $+0.03\%$ \\
 & \textbf{Backslopes ($\theta_{\text{align}} < -0.3$)} & \textbf{5,222} & $\mathbf{53.42\%}$ & $\mathbf{52.36\%}$ & $48.33\%$ & $38.41\%$ & $\mathbf{+5.09\%}$ \\
\cmidrule(lr){2-8}
 & Runout \& Valleys ($< 25^\circ$) & 6,501 & $45.97\%$ & $38.49\%$ & $38.22\%$ & $28.40\%$ & $+7.75\%$ \\
 & Avalanche Chutes ($25^\circ$--$45^\circ$) & 15,875 & $55.82\%$ & $52.33\%$ & $56.25\%$ & $50.55\%$ & $-0.44\%$ \\
 & Extreme Steep ($> 45^\circ$, Shadow) & 203 & $8.85\%$ & $5.04\%$ & $16.65\%$ & $10.74\%$ & $-7.80\%$ \\
\bottomrule
\end{tabular*}
\end{table*}

The Seed-42 empirical stratification provides descriptive, physics-consistent associations:
\begin{enumerate}
    \item \textbf{Backslope Separation in Both Mountain Systems:} Across both held-out mountain regions, TopoRadar-Net achieves its largest Seed-42 margin over Attention U-Net on slopes facing away from the radar look direction ($\theta_{\text{align}} < -0.3$). In Troms{\o}, TopoRadar-Net attains $81.08\%$ F1 versus $75.23\%$ for Attention U-Net, with $+12.28\%$ higher recall (rising from $65.20\%$ to $77.48\%$). This pattern is consistent with the backscatter observations of \citet{tompkin2021backscatter}: on backslopes with high local incidence angles ($55^\circ \pm 20^\circ$), undisturbed snow exhibits low backscatter and rough debris can appear strongly visible. In the high-altitude continental Pamir Mountains, the $+5.09\%$ backslope difference ($53.42\%$ vs $48.33\%$) is precision-associated ($+13.95\%$ precision, rising from $38.41\%$ to $52.36\%$, while recall drops from $65.17\%$ to $54.52\%$). These are descriptive stratum associations from Seed 42, not component-isolated causal effects.
    \item \textbf{Separation in Active Avalanche Chutes ($25^\circ$--$45^\circ$):} In the primary avalanche terrain zone ($25^\circ$--$45^\circ$ slope), TopoRadar-Net attains $81.34\%$ F1 versus $77.00\%$ for Attention U-Net ($+4.34\%$ F1, $+5.94\%$ IoU) in Troms{\o}. The foreslope difference is $-0.07\%$ in this Seed-42 descriptive comparison.
    \item \textbf{Physical Radar Shadow Boundary:} In the extreme high-altitude Pamir Mountains on slopes $> 45^\circ$, performance for both networks is low ($8.85\%$ vs $16.65\%$, $\Delta = -7.80\%$), with debris representing only 203 pixels. The result identifies a small-sample failure stratum consistent with radar shadow; it does not prove that shadow alone causes the errors.
\end{enumerate}

To test whether the exploratory Seed-42 patterns persist under training uncertainty, we repeated full-scene stratification for Seeds 42, 123, and 456 using each checkpoint's frozen validation-selected threshold (multi-seed stratification artifact). The backslope F1 differences versus Attention U-Net average $-0.34\% \pm 6.28\%$ in Troms{\o} and $+1.20\% \pm 7.73\%$ in Pamir (population SD; positive in two of three seeds for each region), rather than stable $>+5\%$ effects. Troms{\o} foreslopes average $-1.21\% \pm 0.86\%$ (zero of three positive), while Pamir runout/valley terrain averages $+5.32\% \pm 7.76\%$ (two of three positive). Extreme-steep strata are highly seed-variable. The multi-seed check therefore narrows Table~\ref{tab:stratification} to an exploratory regional characterization and provides no component-isolated directional mechanism claim.

\subsection{Extended Validation Cohort Analysis on Polar Maritime Terrain (Nuuk, Greenland)}
\label{sec:res:nuuk}

To evaluate model stability across the broader AvalCD distribution beyond the zero-shot test scenes, Table~\ref{tab:nuuk_validation} reports performance on the polar maritime validation cohort in Southwest Greenland (\texttt{Nuuk\_20210411}, used for checkpoint selection and threshold calibration; $n=129$ polygons, 81,595 debris pixels).

\begin{table*}[t]
\caption{Performance comparison on the polar maritime validation scene in Southwest Greenland (\texttt{Nuuk\_20210411}, used for checkpoint selection and threshold calibration; $n=129$ verified avalanche polygons, 81,595 debris pixels) across 3 independent training seeds (Mean $\pm$ Standard Deviation).}
\label{tab:nuuk_validation}
\centering
\footnotesize
\setlength{\tabcolsep}{8pt}
\begin{tabular*}{\textwidth}{@{\extracolsep{\fill}}l@{\hspace{8pt}}c@{\hspace{8pt}}c@{\hspace{8pt}}c@{\hspace{8pt}}c@{}}
\toprule
\textbf{Method} & \textbf{Pixel F1 (\%)} & \textbf{IoU (\%)} & \textbf{Precision (\%)} & \textbf{Recall (\%)} \\
\midrule
SiamUNet-diff & $59.35 \pm 1.58\%$ & $42.22 \pm 1.45\%$ & $58.10 \pm 0.58\%$ & $60.79 \pm 3.70\%$ \\
SiamUNet-conc & $66.53 \pm 0.97\%$ & $49.85 \pm 1.09\%$ & $57.72 \pm 2.15\%$ & $78.66 \pm 1.74\%$ \\
ResU-Net & $67.16 \pm 1.62\%$ & $50.57 \pm 1.81\%$ & $\mathbf{60.10 \pm 4.51\%}$ & $76.69 \pm 3.48\%$ \\
Attention U-Net & $67.01 \pm 1.98\%$ & $50.43 \pm 2.22\%$ & $57.06 \pm 2.94\%$ & $81.31 \pm 0.51\%$ \\
\midrule
\textbf{TopoRadar-Net (Ours)} & $\mathbf{66.75 \pm 1.96\%}$ & $\mathbf{50.13 \pm 2.22\%}$ & $56.63 \pm 3.31\%$ & $\mathbf{81.50 \pm 1.10\%}$ \\
\bottomrule
\end{tabular*}
\end{table*}

In this polar maritime alpine regime, characterized by open fjord expanses and rounded permafrost massifs without extreme knife-edge ridge shadows, all four modern deep change detection architectures cluster within $0.63\%$ of each other on F1 ($66.53\%$--$67.16\%$). TopoRadar-Net achieves the highest detection sensitivity of all models ($81.50 \pm 1.10\%$ recall), successfully capturing the vast majority of the 129 dispersed avalanche tracks, while Attention U-Net ($67.01\%$) and ResU-Net ($67.16\%$) achieve marginally higher precision due to fewer false alarms on open plateau terrain. All deep architectures decisively outperform unguided differencing ($59.35\%$, yielding a $+7.40\%$ to $+7.81\%$ F1 margin). This empirical finding corroborates the regime-dependent sensitivity documented in Section~\ref{sec:res:ablation}: in open sub-polar terrain with moderate relief, unguided baseline models and physics-conditioned networks perform with statistical parity, whereas rigid topographic priors provide their most decisive advantages in steep, fjord-adjacent coastal mountains (Troms{\o}) and high-relief continental massifs (Pamir).
